\section{Determinant estimates by twisting}
\label{sec:determinant-estimates}

The Archimedean estimates concern the Gram determinant for the trivial
twist. Twisting the chain allows us to compare this determinant with
another Gram determinant at each finite place. The choice of twist
controls both the denominators of projections onto spaces of theta lifts
and the Hecke \(L\)-values that enter the comparison. These two bounds
give a determinant ratio at each stage and then a ratio for the full chain.

\subsection{Choosing twists and clearing projection denominators}

Set \(r=\dim_E Y_0\) and \(t=5r\), and write \(W_j,V_j\) for the
Hermitian spaces of ranks \(r,r+t\), respectively, used in the comparison
at stage \(j\). Let \(F_j\) be the image of \(e_0\) on the forms of
principal level on \(W_j\). After projection, the multiplicity of each
Hecke eigensystem at split places away from the level is the same on
\(W_j\) and \(Y_0\). In particular, \(\dim_{\mathbb C} F_j=h\).
\leansrc{New/Spanning/Multiplicities/EqualMultiplicities}{Equality of multiplicities}.

A twist \(\xi\) is an automorphic character of \(U(1)\) of weight \(0\) at
the infinite places. Write \(Q_\xi\) for its set of ramified places.
For every twist under consideration, \(\xi^j\) is ramified at every
place of \(Q_\xi\) for \(1\le j\le m\). The modulus norm appearing
at the level is therefore \(N_\xi=\prod_{w\in Q_\xi}\mathrm Nw\) for
every \(\xi^j\). The twist \(\xi\) is admissible at \(\ell\) if \(N_\xi\)
is coprime to \(\ell\).

Consider a character \(\nu\) of weight \(-c(\tau)\) and an integer \(u\)
with \(1\le u\le t\) and \(u\le c(\tau)\le4b\). They form a critical
pair if the tame part of \(\nu\) is trivial away from the chosen level
places. A critical pair is called ramified if its tame part is nontrivial
at every chosen level place. We use the ratio
\[
 \mathcal L_{\nu,u}(\vartheta;\varepsilon)
 =
 \frac{|L(u,(\nu\vartheta^\varepsilon)^{-1})|}
      {|L(u,\nu^{-1})|},
 \qquad \varepsilon\in\{1,-1\}.
\]
Here \(L\) is a Hecke \(L\)-function over \(E\) with character pulled back
from \(U(1)\) along \(z\mapsto z/\bar z\). Both numerator and denominator
omit Euler factors at the chosen level places and their conjugates.
The numerator also omits the factors at \(Q_\vartheta\) and its conjugate.
A main pair for \(j\) is a critical pair with \(m<u\le r\), with tame
part equal to that of the character at stage \(j\), and with
\(w_j(\tau)+r-t\le c(\tau)\le w_j(\tau)\) at the infinite places, where
\(-w_j(\tau)\) is the weight of that stage's character. For a main pair,
we decompose the associated space of rank \(2u\) into four blocks of
ranks \(12,4,4,2u-20\).
\leansrc{New/Determinant/Twist/FitTwists}{Critical pairs and main pairs}.

Write \(u_{j,\vartheta}\) for the theta lift on \(F_j\) and
\(P_{j,\vartheta}\) for the orthogonal projection onto its image
\(I_{j,\vartheta}\). At a split place \(w\) away from the level, specialize
the Hecke polynomial operator at the value associated with the lift at
stage \(j\). The false lifts at \(w\) form the intersection of the kernel
of this operator, the image of \(e_0\), and \(I_{j,\vartheta}^{\perp}\).
Orthogonality to false lifts means orthogonality at every such \(w\).
\leansrc{New/OneLift/TestedForms}{Definition of false lifts}.

Put
\[
 B_{\mathrm{den}}=\frac{E_3}{mk}.
\]
The existence theorem for regular twists supplies, for each \(\iota\),
a twist \(\xi_\iota\) admissible at \(\ell\), with the three properties
below. As before, we write \(|v|_\iota=|\iota(v)|_\ell\).
\begin{enumerate}
\item For every ramified pair, every \(1\le j\le m\), and every
  \(\varepsilon=\pm1\), the family
  \(\bigl(\mathcal L_{\nu,u}(\xi_\iota^j;\varepsilon)^{-1}\bigr)_\iota\)
  has height at most \(x_1\).
\item For each \(j\), there are nonnegative reals \(\beta_\ell\) indexed
  by primes, with \(\sum_{\ell\in T}\beta_\ell\le B_{\mathrm{den}}\) for
  every finite set \(T\) of primes. For \(\vartheta=\xi_\iota^j\), there
  is an endomorphism \(\Lambda\) of \(I_{j,\vartheta}\) such that
  \[
    |\det\Lambda|_\iota\ge\exp(-hr\beta_\ell).
  \]
  The operator \(\Lambda P_{j,\vartheta}\) takes forms that are integral
  at \(\iota\) and orthogonal to the false lifts to forms integral
  at \(\iota\).
\item The same \(\beta_\ell\) controls projection onto the line generated
  by the lift from rank \(1\) associated with a main pair. Write \(P\)
  for the orthogonal projection onto that line. If \(\Xi,\Xi'\) are
  test forms integral at all places above \(\ell\) and orthogonal to
  the false lifts for the rank-one descent, then
  \[
       |\langle P\Xi,P\Xi'\rangle|_\iota\le\exp(\beta_\ell).
  \]
\end{enumerate}

The existence theorem uses the Galois representations of rank \(r\)
attached to the finitely many Hecke eigensystems in the \(e_0\)-part of
the forms on \(W_m\), together with the rank-one representation attached
to each main pair. The twisting theorem is applied to the main summands
whose traces are obtained from those of the twisted lifts by removing a
string of characters. For the first \(m\) powers of one anticyclotomic
character it bounds both the lengths of the Selmer groups with divisible
coefficients and the valuations of the \(L\)-value ratios for ramified pairs.
Let \(C_{\mathrm{lev}}\) be the product of level group orders used in the
formal proof, and \(C_1\) the number of classes of \(U(1)\) at the chosen
principal level. For the inverse ratio of a ramified pair the bound for
\(\log|\cdot|_\iota\) is
\[
  \big((4f+3)r\,v_\ell(C_{\mathrm{lev}})+2^{-\ell}\big)\log\ell.
\]
The denominator estimate obtained by Ribet's method combines the Selmer
group length with local and level corrections in the exponent \(c_\ell\).
\leanref{clearingBound}
The inner product of rank-one projections also requires an allowance of
\(v_\ell(C_1)\log\ell\), so both estimates use
\(\beta_\ell=(c_\ell+v_\ell(C_1))\log\ell\).
The two sums over primes are bounded by \(x_1\) and \(B_{\mathrm{den}}\),
respectively. We use the twisting theorem and the passage from Selmer
group lengths to projection denominators in this form.
\leansrc{New/Congruences/Selmer/TwistingTheorem}{Twisting theorem for Selmer groups},
\leansrc{New/Congruences/FirstOrder/DenominatorBound}{Projection denominator bound},
\leansrc{New/Congruences/FitTwistsExist}{Existence of regular families}.

\subsection{Hecke values and theta derivatives for main pairs}

Fix a main pair \(p=(\nu,u)\) and write \(\phi_\vartheta\) for the normalized
lift from rank \(1\) to rank \(2u\) associated with it. Comparing the norm
formula for the character lift at the two twists gives
\begin{equation}
 \mathcal L_p(\vartheta;1)^{-1}
   =N_\vartheta^{2u-1}\frac{\|\phi_1\|^2}{\|\phi_\vartheta\|^2}.
 \label{eq:main-pair-norm}
\end{equation}
To estimate the ratio at finite places, we project two test forms onto the
line spanned by \(\phi_\vartheta\).

First, consider the height of \(\|\phi_1\|^2\) as a constant family.
Let \(\phi^{(0)}_1\) be the lift corresponding to a critical pair of the
same rank and derivative order \(0\). Its squared norm becomes integral
at every finite place after multiplication by an integer \(n_L\).
Choose \(n_L\|\phi^{(0)}_1\|^2/\|\phi_1\|^2\) as the common multiplier
in the height bound for \(\|\phi_1\|^2\). This makes the finite-place
contribution zero. At infinity, the two logarithmic squared norms have
a common term that cancels. The derivative order for the main pair is
\(c(\tau)-u\ge\kappa(\tau)\), so the contribution from the derivative
orders differs from that for order \(0\) by at least
\(\sum_\tau s_\infty(\kappa(\tau))\). The error in the norm formula
together with \(\log n_L\) is bounded by \(B_{\mathrm{den}}\), giving
\begin{equation}
 \operatorname{ht}\big((\|\phi_1\|^2)_\iota\big)
  \le B_{\mathrm{den}}-\sum_{\tau} s_\infty(\kappa(\tau)).
 \label{eq:untwisted-rankone-height}
\end{equation}
The norm estimate is first proved for conjugations fixing the quadratic
field \(K\). The lemma on absolute values under conjugation extends it
to every field automorphism of \(\mathbb C\).
\leansrc{New/OneLift/HeckeValues/AlgebraicPartSize}{Height of the untwisted norm}.

Write the required finite-place valuation bound for the theta derivative as
\[
 F(\ell)=\frac12\sum_{\tau\in\Phi}\sum_{e\ge 1}
     (1-\chi_K(\ell^e))
     \left(\frac{\kappa(\tau)}{\ell^e}-\kappa(-\ell^e\tau)\right)_+,
\]
where \(\chi_K=\chi_d\) is the quadratic character of \(K\).
Grouping the terms by prime powers gives, for every finite set \(T\)
of primes,
\begin{equation}
 2\sum_{\ell\in T}F(\ell)\log\ell
   \le \sum_{\tau\in\Phi}s_{\mathrm{fin}}(\kappa(\tau),\tau).
 \label{eq:local-prime-total-lower}
\end{equation}
\leansrc{New/UpperBound/Size/Estimates}{Sum of local prime contributions}.

The construction of the test form expands the lift of rank \(2u\)
according to the four blocks above. In the block of rank \(12\), choose a
theta-series coefficient that attains the valuation bound. This supplies
a term of valuation at most \(F(\ell)\) plus the valuation of a fixed
factor. In the second block, of rank \(4\), use the derivative-order-zero
lift of a ramified pair \(q_2\), with ratio
\(\rho_\vartheta=\mathcal L_{q_2}(\vartheta;1)\). The remaining two blocks
are assigned the two factors of a character of the class group of
\(U(1)\). Translate the class representatives so that the trivial class
is represented by the identity. Summing the normalized, character-weighted
class sums over all characters recovers the chosen coefficient by character
orthogonality. The ultrametric inequality then gives at least one class
sum whose \(\iota\)-absolute value is at least that of the coefficient.
Choose a factorization of the corresponding character so that the product
of the algebraic \(L\)-values from the last two blocks has
\(\iota\)-absolute value at least their geometric mean over all possible
factorizations. The product over the blocks defines the test form.
Expanding its pairing with the lift as a class sum yields these factors.

The resulting local estimate uses a scalar
\(c_0\ne0\), independent of the twist and the finite place, with
\(|\sigma(c_0)|\le\exp(fmb/k+3x^{34})\le\exp(B_{\mathrm{den}})\) for every
\(\sigma\). For each \(\iota\) and each \(\vartheta\) admissible at its
prime, there is a test form \(\Xi\) such that
\begin{equation}
 |\langle\Xi,\phi_\vartheta\rangle|_\iota
    \ge |c_0\rho_\vartheta|_\iota\,\ell^{-F(\ell)}.
 \label{eq:test-pairing-lower}
\end{equation}
Here \(c_0\) is the product of the fixed factor multiplying the theta
coefficient, the untwisted algebraic \(L\)-value from the second block,
and the geometric mean from the remaining two blocks. The same construction
makes the test form integral at all places above \(\ell\) and orthogonal
to the false lifts.
\leansrc{New/Derivatives/Sharpness/BoundAttained}{Theta coefficient valuations},
\leansrc{New/OneLift/HeckeValues/TwistedValueHeights}{Construction of the test form}.

Take \(\vartheta=\xi_\iota^j\). Choose test forms \(\Xi,\Xi'\)
as in \eqref{eq:test-pairing-lower} at \(\iota\) and
\(\iota\circ\mathrm{conj}\), respectively. Since \(\rho_\vartheta\)
is real, the product
\(A_\iota=\langle\Xi,\phi_\vartheta\rangle\langle\phi_\vartheta,\Xi'\rangle\)
satisfies
\[
 \left|c_0\bar c_0\,\frac{\rho_\vartheta^2}{A_\iota}\right|_\iota
      \le \ell^{2F(\ell)}.
\]
Both pairings are nonzero by \eqref{eq:test-pairing-lower}.
The orthogonal projection \(P\) onto \(\mathbb C\phi_\vartheta\) gives
\(\langle P\Xi,P\Xi'\rangle=A_\iota/\|\phi_\vartheta\|^2\).
Thus the right-hand side of \eqref{eq:main-pair-norm} is \(N_\vartheta^{2u-1}\)
times the product of four factors
\[
 \|\phi_1\|^2,\qquad
 \langle P\Xi,P\Xi'\rangle,\qquad
 \rho_\vartheta^{-2},\qquad
 \frac{\rho_\vartheta^2}{A_\iota}.
\]
The height bound for each factor follows, in order, from
\eqref{eq:untwisted-rankone-height}, the estimate for the projected
pairing, the inverse-ratio bound for ramified pairs, and
\eqref{eq:local-prime-total-lower} with common multiplier \(c_0\bar c_0\).
The four bounds are
\[
 \begin{gathered}
 B_{\mathrm{den}}-\sum_\tau s_\infty(\kappa(\tau)),\qquad
 B_{\mathrm{den}},\qquad 2x_1,\\
 2B_{\mathrm{den}}+\sum_\tau s_{\mathrm{fin}}(\kappa(\tau),\tau).
 \end{gathered}
\]
Moreover, \(N_\vartheta^{2u-1}\) is a unit at \(\iota\).
Adding the bounds and using \(x_1\le B_{\mathrm{den}}\) gives
\begin{equation}
 \operatorname{ht}\bigl(
   (\mathcal L_p(\xi_\iota^j;1)^{-1})_\iota\bigr)
 \le 6B_{\mathrm{den}}-A_0.
 \label{eq:main-pair-height-lower}
\end{equation}
\leansrc{New/OneLift/HeckeValues/TwistedValueHeights}{Height of the inverse ratio for a main pair}.

\subsection{Comparing determinants at a single stage}

Set \(Z_j(\vartheta)=\det(u_{j,\vartheta}^*u_{j,\vartheta})\). The
comparison operator \(U_{j,\vartheta}\) is a scalar multiple of
\(G_{j,\vartheta}G_{j,\vartheta}^*\), where \(G_{j,\vartheta}\) glues
tuples of forms on an auxiliary space into forms on \(V_j\). The scalar
is an integer from the level times the square of a hook number.
Three exchanges of theta integrals give a positive-definite operator
\(U_j^\circ\) on \(F_j\), independent of \(\vartheta\), such that
\begin{equation}
 u_{j,\vartheta}^*U_{j,\vartheta}u_{j,\vartheta}
      =\rho_j(\vartheta)U_j^\circ,\qquad
 \rho_j(\vartheta)=
 \frac{\mathcal H_j(\vartheta)}{\mathcal H_j(1)}.
 \label{eq:comparison-operator-lower}
\end{equation}
Here \(\mathcal H_j\) is the product of the three Hecke values that enter
the exchanges. \leanref{heckeValueProd}
The operator \(U_j^\circ\) is a positive scalar multiple of the Gram
operator of the untwisted period map. This map is injective, so
\(\det U_j^\circ\ne0\).

The ratio \(\rho_j(\vartheta)\) is nonzero. The operator \(U_{j,\vartheta}\)
carries theta lifts into forms orthogonal to the false lifts and preserves
integrality at primes coprime to \(N_\vartheta\). For the second property,
a tight frame of integral test vectors expresses
\(G_{j,\vartheta}G_{j,\vartheta}^*\) as a sum of rank-one operators.
The level integer clears the denominators of their pairings.
The Cauchy--Schwarz estimate for the same frame shows that the untwisted
operator increases the norm by at most \(\exp(3fmb(\log b)^2)\) at every
conjugate.
\leansrc{New/OneLift/ComparisonOperator/ThreeExchanges}{The three exchanges},
\leansrc{New/OneLift/LiftVolumeHeight}{The comparison operator and its bounds}.

Let \(B_{j,\vartheta}\) be the restriction of
\(P_{j,\vartheta}U_{j,\vartheta}\) to \(I_{j,\vartheta}\) and put
\(D_j(\vartheta)=\det B_{j,\vartheta}\). The lift \(u_{j,\vartheta}\)
is an isomorphism onto its image, and \(u_{j,\vartheta}^*\) vanishes
on \(I_{j,\vartheta}^\perp\). Hence
\[
 u_{j,\vartheta}^* U_{j,\vartheta} u_{j,\vartheta}
  =(u_{j,\vartheta}^*u_{j,\vartheta})
    (u_{j,\vartheta}^{-1} B_{j,\vartheta}u_{j,\vartheta}).
\]
Taking determinants and using \eqref{eq:comparison-operator-lower} gives
\[
 Z_j(\vartheta)D_j(\vartheta)
      =\det(U_j^\circ)\rho_j(\vartheta)^h,\qquad
 \frac{Z_j(1)}{Z_j(\vartheta)}
   =\frac{D_j(\vartheta)}{D_j(1)}\,\rho_j(\vartheta)^{-h}.
\]

At finite places, for \(\vartheta=\xi_\iota^j\) the operator
\(\Lambda B_{j,\vartheta}\) preserves the integral forms in
\(I_{j,\vartheta}\). Normalize each eigenvector so that it is integral
and rescaling it by a scalar with \(\iota\)-absolute value greater than
\(1\) takes it outside the integral module. The corresponding eigenvalue
has \(\iota\)-absolute value at most \(1\). Consequently,
\[
 |D_j(\xi_\iota^j)|_\iota
   \le|\det\Lambda|_\iota^{-1}
   \le\exp(hr\beta_\ell).
\]
At infinity, \(U_{j,1}(I_{j,1})\) is orthogonal to the false lifts.
Projecting forms in this image back onto \(I_{j,1}\) cannot increase
their norm at any conjugate. The projection estimate uses the adjoint
relation for the specializations of the Hecke operators under conjugation.
Put \(C_{\mathrm{arch}}=fmb(\log b)^2\). Multiplying the bounds for the
eigenvalues gives \(|\sigma(D_j(1))|\le\exp(3h C_{\mathrm{arch}})\).
Using \(D_j(1)\) as common multiplier gives
\[
 \operatorname{ht}\bigl((D_j(\xi_\iota^j)/D_j(1))_\iota\bigr)
     \le h(3C_{\mathrm{arch}}+rB_{\mathrm{den}}).
\]
\leansrc{New/OneLift/LiftVolumeHeight}{Finite and Archimedean bounds for the compressed determinant}.

Up to an integer power of \(N_\vartheta\), the ratio \(\rho_j(\vartheta)\)
in the three exchanges is the square of the product of the ratios for
\(k/2\) main pairs and \(5m\) ramified pairs. The main pairs occur at points
\(7m/2<u\le 7m/2+k/2\). The remaining points number \(7m/2+m+m/2=5m\).
The sign is \(\varepsilon=1\) for the main pairs and for the first \(7m/2\)
ramified pairs. Among the flat pairs it is \(\varepsilon=1\) for \(m\) of
them and \(\varepsilon=-1\) for \(m/2\) of them.
Adding the bounds in \eqref{eq:main-pair-height-lower} and those for the
ramified pairs, with the factor of two from the square, gives
\[
 \operatorname{ht}\bigl((\rho_j(\xi_\iota^j)^{-1})_\iota\bigr)
  \le k(6B_{\mathrm{den}}-A_0)+10mx_1.
\]
The admissible data satisfy
\(3C_{\mathrm{arch}}+(r+6k)B_{\mathrm{den}}+10mx_1
        \le kB_{\mathrm{den}}\log b\), so for every stage
\begin{equation}
 \operatorname{ht}\bigl((Z_j(1)/Z_j(\xi_\iota^j))_\iota\bigr)
   \le h\big(E_4/m-kA_0\big).
 \label{eq:single-relative-height}
\end{equation}
\leansrc{New/OneLift/HeckeValues/TwistedValueHeights}{Factorization of Hecke values},
\leansrc{New/OneLift/LiftVolumeHeight}{Height bound for a single stage}.

\subsection{The determinant of the chain}

To compute the norm of the chain, arrange in a triangle the lifts that
adjoin blocks of height \(t\) to the initial space. The vertex \((i,n)\)
consists of the common part and blocks \(i+1,\ldots,i+n\). The arrows
that add a block on the right are paired with transposed lifts that
remove a block on the left. Comparing the two paths obtained by switching
the order of these operations uses the exchange identity.

The matching over a finite ring works as follows. Split the row and column
indices into two groups. For an invertible matrix \(g\), write \(C(g)\)
for its first columns and \(R(g)\) for the last rows of \(g^{-1}\), with
the sizes specified by the split. The weights \(\Psi,\Phi\) are supported
on matrices of these respective forms and satisfy
\(\Phi(R(g))=\chi(\det g)\Psi(C(g))\).
Let \(K,K'\) be the kernels along the two paths, subject to the formal
proof's weight, ramification, signature, and rank conditions.
The exchange gives \(\sum\Psi K'=(\kappa_0/c)\sum\Phi K\).
The value \(|c|^2\) is determined by the product of powers of residue-field
cardinalities and squared absolute values of partial \(L\)-values.
The scalar \(\kappa_0\) is independent of the twist.
Applying the identity to point data and to their finite Fourier transforms
gives the two exchange laws needed for the triangle.
\leansrc{New/Determinant/Exchange/TheoremForPartners}{Exchange of theta integrals},
\leansrc{New/Determinant/Exchange/OneSquare}{Exchange across one square},
\leansrc{New/Determinant/Chain/ChainNorms}{Exchange laws and injectivity for the triangle}.

Write \(c_Q(\xi)\) for
\[
 N_\xi^{rt}\prod_{w\in Q_\xi}\prod_{t<s\le t+r}(1-(\mathrm Nw)^{-s}).
\]
Let \(R_j(\xi^j)\) be the Hecke operator on \(V\) from Rallis's formula
at stage \(j\). Applying Rallis's formula to a single lift and matching
the multiplicities of the Hecke eigensystems of \(W_j\) and \(Y_0\)
gives, for a constant \(A_j\ne0\) independent of the twist,
\[
 Z_j(\xi^j)=(A_jc_Q(\xi))^h\det R_j(\xi^j).
\]
For the chain, the exchange laws move the blocks one at a time.
Write \(X_{i,n}\) for the composition of the first \(n\) arrows in row
\(i\), weighting the arrow from stage \(s\) to \(s+1\) by the order
\(w_s\) of the general linear group at its source's finite level.
Let \(a_{i,n}\) be the arrow to the right and \(d_{i,n}\) the downward
arrow to the next row. The map \(j_i\) embeds the space of principal-level
forms isometrically at the beginning of row \(i\). Define \(x'\) by
\(j_{i+1}x'=d_{i,0}a_{i,0}j_i x\).
Composing the exchange laws for individual squares and using the
exchange law for adjoints gives
\begin{multline*}
  \langle X_{i,n+1}j_i x,X_{i,n+1}j_i y\rangle\\
   = \lambda_{i,n}\,\frac{w_0w_1w_n}{w_{n+1}}\,
     \langle X_{i+1,n}j_{i+1}x',X_{i+1,n}j_{i+1}y\rangle.
\end{multline*}
Up to a factor independent of the twist and a half-integral power of
\(N_\xi\), the absolute value of \(\lambda_{i,n}\) is
\(\prod_{i+1<j\le i+n+1}J_{i+1,j}(\xi)\).
On the right, the pairing with \(x'\) becomes a pairing with \(x\)
by applying the Gram operator of the first arrow in row \(i\),
canceling one \(w_0\). Starting the induction with the empty
composition, the Gram operator for row \(i\) is a scalar multiple of the
product of the Gram operators of the first arrows in the successive rows.
The scalar includes \(w_0w_1^n/w_n\) and the contributions from pairs
\(i<r'<j\le m\).
\leansrc{New/Determinant/Chain/TriangleNorms}{Recurrence for the weighted chain}.

For the chain starting in row \(0\), apply Rallis's formula to the
Gram operator of the first arrow in each row. Expressing the final level
index by a count of block upper triangular matrices cancels
\(w_0w_1^m/w_m\). With this normalization the Gram operator is
\[
 \Theta_\xi^*\Theta_\xi
       = c(\xi)c_Q(\xi)^m \prod_{j=1}^m R_j(\xi^j),
\]
where
\[
 |c(\xi)| = C\,N_\xi^{e/2}\prod_{1\le i<j\le m} J_{ij}(\xi),
 \qquad C>0,\quad e\in\mathbb Z.
\]
Write \(\eta_j=\mu_j^0\) for the untwisted character at stage \(j\).
Then
\[
 J_{ij}(\xi)=
   \prod_{u=1}^{t}\mathcal L_{\eta_j/\eta_i,u}(\xi^{j-i};1).
\]
Every pair \(i<j\) arises in moving the blocks. The tame part of
\(\eta_j/\eta_i\) at the chosen level is the \((j-i)\)-th power of a
fixed character. Since \(1\le j-i\le m\), the pairs in this product are
ramified.

Injectivity of the chain follows from that of the weighted chain in the
triangle. The exchange laws imply that when the weighted chain of length
\(n+1\) kills \(j_i v\), the chain of length \(n\) in the next row kills
\(d_{i,0}a_{i,0}j_i v\). This vector lies in the image of \(j_{i+1}\).
Induction on the length and injectivity of \(d_{i,0}a_{i,0}\) on the
image of \(j_i\) give \(v=0\). The positive level index and the positive
dimension of the Weyl module that enter the change of normalization
transfer this injectivity to \(\Theta_\xi\). Each individual lift is
also injective, so all the Gram determinants are positive. Taking
absolute values on both sides of the determinant identities retains
these positive Gram determinants and cancels \(c_Q(\xi)\) and the
determinants of the Hecke operators. Thus
\begin{equation}
 \frac{Z(1)}{Z(\xi)}
 =
 \left(N_\xi^{e/2}\prod_{1\le i<j\le m}J_{ij}(\xi)\right)^{-h}
       \prod_{j=1}^m\frac{Z_j(1)}{Z_j(\xi^j)}.
 \label{eq:chain-relative-factorization}
\end{equation}
If \(\ell\nmid N_\xi\), then \(N_\xi^{e/2}\) has \(\iota\)-absolute value
\(1\). Counting the ramified pairs gives height at most \(m^2t x_1\) for
the inverse of the parenthesized factor evaluated at \(\xi_\iota\), and
\(m^2t x_1\le b\). Adding \eqref{eq:single-relative-height} over the
\(m\) stages now gives
\begin{equation}
 \operatorname{ht}\bigl((Z(1)/Z(\xi_\iota))_\iota\bigr)
       \le h(b+E_4-H_0).
 \label{eq:chain-relative-height}
\end{equation}
\leansrc{New/Determinant/Chain/NormOfChain}{The chain's Gram operator and determinant factorization}.

\subsection{Forms for completing a basis}

To complete a basis, we span the space of principal-level forms before
applying \(e_0\) using forms of deep level. The formal proof uses two
grids. It transfers the period of the initial path along their boundary
to the period of the final path, making one turn in each of the \(m^2\)
squares. At each turn it applies the exchange identity to a sum over
finite matrix data. At a turn on the diagonal it first identifies the
relation to a rank-one kernel and then applies the same comparison.
The resulting nonzero scalar \(c_{\mathrm{path}}(\xi)\) satisfies
\begin{equation}
 \left|\frac{c_{\mathrm{path}}(\xi)}{c_{\mathrm{path}}(1)}\right|
   =N_\xi^{e'/2}\prod_{0\le i,j<m} T_{ij}(\xi),\qquad
 T_{ij}(\xi)=
    \prod_{u=1}^{t}\mathcal L_{\eta_{i+1}\eta_{j+1},u}(\xi^{j+1};1).
 \label{eq:turn-ratio-lower}
\end{equation}
The exponent of the tame character is \(i+j+2\), between \(2\) and \(2m\),
so the tame part is nontrivial at every chosen level place. Thus the
pairs in each \(T_{ij}\) are ramified.
\leansrc{New/Determinant/AuxiliaryForms/Turns}{The \(L\)-value ratio for a turn},
\leansrc{New/Determinant/AuxiliaryForms/DiagonalTurns}{Comparison of the full paths}.

The final path period sums the inner product of the lifts along the two
half-paths over tuples \(S\) of finite matrices, with coefficient \(\Psi(S)\).
With the normalization fixed, an integer independent of \(\xi\) clears the
denominators of the kernels along the edges and of the class sums when
\(\ell\nmid N_\xi\). For the trivial twist, the Weyl norm bounds for the
kernels and the number of summands give a bound for every conjugate of
the period. To recover the values of the original lift, pair it with dual
delta forms. For the matrix \(D\) of the dual pairing there are an integral
matrix \(C\) and a positive hook number \(H\) with \(CD=HI\).
If the pairing against each integral vector lies in \(H\) times the
valuation ring, testing against the rows of \(C\) and canceling \(H\) shows
that each coordinate of the original value is integral. At infinity, dividing
the bound for the pairing by the norm of the dual vector gives the Weyl
norm bound.

This gives a quantitative spanning result for the principal-level space.
There is a surjection \(s\) from the deep-level space to the principal-level
space and a constant \(c_{\mathrm{sc}}\ne0\). The deep-level space is
spanned by forms \(\psi\) integral at every finite place and with all
conjugates of their values of squared Weyl norm at most \(\exp(x^{99}b^2)\).
For each \(\psi\), scaling the lift of \(c_{\mathrm{sc}}s(\psi)\) by
\(c_{\mathrm{path}}(\xi)/c_{\mathrm{path}}(1)\) makes it integral at
admissible places. For the trivial twist, the squared Weyl norms of the
lift's values are bounded by the deep-level bound times
\(\exp(2x^{201}b^2)\), hence by \(\exp(x^{202}b^2)\).
These are the local and Archimedean kernel estimates used in the
construction.
\leansrc{New/Determinant/AuxiliaryForms/SpanningForms}{The final path period and construction of auxiliary forms}.

Send this spanning family to \(V\) by the projection multiplied by
the order of the level group. The scaled projection preserves
integrality and adds at most \(x^{99}\) to the logarithm of the
squared norm for the trivial twist. Since
\(x^{99}+x^{202}b^2\le2x_1^2b^2\), the resulting auxiliary forms \(z\)
satisfy
\[
 \|\Theta_1 z\|_\sigma^2\le \exp(2x_1^2b^2)
 \qquad(\sigma\in\operatorname{Aut}(\mathbb C)).
\]
Set \(\gamma(\xi)=c_{\mathrm{path}}(\xi)/c_{\mathrm{path}}(1)\).
Then \(\gamma(\xi)\Theta_\xi z\) is integral at every admissible place.
Equation \eqref{eq:turn-ratio-lower} and the bounds for ramified pairs give
\begin{equation}
 \operatorname{ht}\bigl((|\gamma(\xi_\iota)|^{-2})_\iota\bigr)
 \le 2m^2t x_1\le b.
 \label{eq:auxiliary-factor-height}
\end{equation}
\leansrc{New/Determinant/Twist/LongForms}{Bounds for auxiliary forms after projection to the target}.

\begin{proposition}\label{prop:determinant-inputs}
For all admissible data, \(\Theta_\xi\) is injective for each twist \(\xi\).
There exist a twist \(\xi_\iota\) admissible at the prime of each \(\iota\)
and constants \(\gamma(\xi)\ne0\), with \(\gamma(1)=1\), such that
\[
 \operatorname{ht}\bigl((Z(1)/Z(\xi_\iota))_\iota\bigr)
       \le h(b+E_4-H_0),\qquad
 \operatorname{ht}\bigl((|\gamma(\xi_\iota)|^{-2})_\iota\bigr)\le b.
\]
The space \(V\) is spanned by auxiliary forms \(z\) for which
\(\gamma(\xi)\Theta_\xi z\) is integral at every admissible place and
\(\|\Theta_1 z\|_\sigma^2\le\exp(2x_1^2b^2)\) for every \(\sigma\).
\end{proposition}

\begin{proof}
The comparison of lifts in the triangle gives injectivity.
For a regular family, \eqref{eq:single-relative-height} holds at
each stage. Substituting those bounds into
\eqref{eq:chain-relative-factorization} gives
\eqref{eq:chain-relative-height}. The surjection from deep level and
the projection by \(e_0\) give the spanning family of auxiliary forms,
with \eqref{eq:auxiliary-factor-height} for the same scalars \(\gamma(\xi)\).
\end{proof}
